% =====================================================================
%  21.3  Chairman / prefix-version trees: the k-th smallest in a range
%  Exports: sec:c19-chairman  def:c19-roots  thm:c19-kth
%           lst:c19-chairman  lst:c19-kth  lst:c19-mkthnum-soln
%           prob:c19-mkthnum  fig:c19-roots
%  All numbers agree with 00_opener and the running array:
%    A = (3,2,4,5,1,1,5,3,7,0,2,4,1,6,3,2), sum 49, 16 slots.
%  Chairman traces (frequency over the VALUE axis 0..7), verified by hand:
%    prefix histograms F_i[value] for values 0,1,2,3,4,5,6,7:
%      F_2 =(0,0,1,1,0,0,0,0)  F_8 =(0,2,1,2,1,2,0,0)  F_9 =(0,2,1,2,1,2,0,1)
%    kth on [3..9] = F_9 - F_2 = (0,2,0,1,1,2,0,1) (multiset 1,1,3,4,5,5,7)
%      k=4: left[0..3]=3 <4 -> go right k'=1; [4..5] count 3 >=1 -> left; [4]=1>=1
%      -> value 4.  kth[3..9](k=4)=4.   (sorted 1,1,3,4,5,5,7 -> 4th = 4.)
%    kth on [1..8] = F_8 - F_0 = (0,2,1,2,1,2,0,0) (multiset 1,1,2,3,3,4,5,5)
%      k=5: left[0..3]=5 >=5 -> left; [0..1]=2 <5 -> right k'=3; [2]=1 <3 ->
%      right k''=2 -> value 3.  kth[1..8](k=5)=3.  (sorted 5th = 3.)
%  SPOJ MKTHNUM (live-verified on spoj.com/problems/MKTHNUM): "K-th Number".
%    n<=100000, m<=5000, n DISTINCT integers |a_i|<=1e9; query i j k asks the
%    k-th number of a[i..j] if that segment were sorted ascending (1<=k<=j-i+1).
%    Statement's own worked example a=(1,5,2,6,3,7,4), Q(2,5,3): segment
%    (5,2,6,3) sorted (2,3,5,6) -> 3rd = 5.
%    Official sample: n=7,m=3, a=(1,5,2,6,3,7,4):
%      (2,5,3)->5; (4,4,1)->6; (1,7,3)->3.  Output: 5 6 3.
% =====================================================================
\section{Chairman Trees: the \(k\)-th Smallest in a Range by Two Root Descents}\label{sec:c19-chairman}

\begin{previouslybox}[The two engines so far each pay a \(\log\) they cannot remove]
Chapter~19 (\chref{c13}) could fold one scalar per interval, and \(\secref{c19-mst}\) kept a whole sorted list per interval to \emph{count} values \(\le x\) -- but at \(\bigO{\log^2 n}\) per count, so its \(k\)-th smallest is a \(\bigO{\log^3 n}\) pipeline (bisect the value, and each test re-tiles the range). \(\secref{c19-copy}\) then taught us to \emph{keep} a whole family of trees, one per version, cheaply, by path copying. This section fuses the two ideas with a choice so natural it feels forced: let the versions be the array's own \term{prefixes}. Version \(i\) is the \emph{frequency histogram of \(A[1..i]\)} stored as a segment tree over the \emph{value} axis; then a range \(A[l..r]\) is the difference of two versions -- subtract \(\mathit{root}_{l-1}\) from \(\mathit{root}_r\) -- and the \(k\)-th smallest drops out of a \emph{single joint descent} down those two trees in \(\bigO{\log V}\). This is the \term{chairman tree} (also called a \term{persistent segment tree}, and historically the \term{functional segment tree}); it is exactly the engine Ch.~19's decision flowchart and Ch.~20's closing lastword both promised would ``do the \(k\)-th smallest inside a static range faster than a merge-sort tree.''
\end{previouslybox}

\begin{topicintro}
Sort the distinct values of \(A\) into \(\mathit{vals}=[v_1<\dots<v_V]\) and let \(\mathit{rank}(x)\) be a value's index in that list (coordinate compression -- the tree lives on ranks \(1..V\), not on the raw values, so huge or negative \(a_i\) cost nothing extra). Build a chain of roots \(\mathit{root}_0,\mathit{root}_1,\dots,\mathit{root}_n\) where \(\mathit{root}_0\) is the empty frequency tree and \(\mathit{root}_i=\) \texttt{update}\((\mathit{root}_{i-1},\,\mathit{rank}(A_i))\): a path copy that walks to leaf \(\mathit{rank}(A_i)\) and increments \(\mathit{cnt}\) on every node of that one path, sharing the other children. So \(\mathit{root}_i\) stores ``how many of \(A[1..i]\) fall in each value sub-range,'' in a \(\bigO{\log V}\)-node-per-insert structure.

The payoff is the \term{two-root subtraction}. For any value sub-range \([p,q]\),
\[
  \#\{\,i\in[l,r]: \mathit{rank}(A_i)\in[p,q]\,\}
  \;=\; \mathit{root}_r[p..q].\mathit{cnt} \;-\; \mathit{root}_{l-1}[p..q].\mathit{cnt},
\]
because \(A[1..r]\) minus \(A[1..l-1]\) is exactly \(A[l..r]\). To get \(\mathit{count}(\le x\ \text{in}\ [l,r])\) we read two node counts and subtract, \(\bigO{\log V}\); to get the \(k\)-th smallest we stand \emph{simultaneously} at \(\mathit{root}_{l-1}\) and \(\mathit{root}_r\) and descend: the difference of their two left-child counts says how many of \(A[l..r]\) lie in the lower value-half, and that single number decides which half \(k\) lives in.

\textbf{Prerequisites.} \(\secref{c19-copy}\) (path copying, the node pool, \(\coderef{c19-pseg}\)) is assumed -- chairman \emph{is} that machinery with versions \(=\) prefixes. Ch.~19 (\chref{c13}) owns the interval/split geometry and the \(\secref{c13-search}\) ``descend the tree vs.\ binary-search outside'' contrast, which we invoke once. \(\secref{c19-ontrees}\) replants these roots on a graph; this section is the flat-array case.

\textbf{After this section you can}
\begin{itemize}[nosep]
  \item build the prefix-version chain and read \(\mathit{count}(\le x\ \text{in}\ [l,r])\) as a two-root subtraction; \hfill\lv{3}
  \item answer the \(k\)-th smallest of \(A[l..r]\) by one joint \(\bigO{\log V}\) descent of two roots; \hfill\lv{4}
  \item apply coordinate compression so raw values up to \(10^9\) cost \(\bigO{\log n}\) nodes; \hfill\lv{3}
  \item recognise that ``static \(k\)-th in a range / count-in-range at a tight budget'' is chairman's job, and that \(\bigO{\log^3 n}\) merge-sort is the slower cousin. \hfill\lv{3}
\end{itemize}
\end{topicintro}

\begin{recognitionbox}[Pattern card -- ``\(k\)-th smallest in a static subarray, or count values in a range, many queries'']
\reckey{many queries ``the \(k\)-th smallest of \(a[l..r]\)'' / ``how many of \(a[l..r]\) are \(\le x\)'' on \emph{one fixed array}; `order statistics on a subarray'; a merge-sort-tree \(\bigO{\log^3 n}\) that just went \tle and needs \(\bigO{\log n}\); counting \emph{distinct} or \emph{frequency} inside a range where the array is read right-to-left as growing prefixes.}
\recuse{
  \item static \(k\)-th smallest in \(a[l..r]\), \(n\le10^5\) and \(\bigO{n\log n}\) memory is fine -- this section's \code{kth} (\(\coderef{c19-kth}\)), SPOJ~MKTHNUM;
  \item range count of a value/threshold answered \emph{online} in \(\bigO{\log n}\) -- the two-root subtraction, replacing an offline BIT sweep;
  \item values up to \(10^9\) or negative -- coordinate-compress to ranks; the tree never sees the raw numbers.
}
\recnotuse{
  \item the array is \emph{updated} between queries -- prefix versions go stale on a write; that is a mutable-tree or a BIT problem (Ch.~19 \chref{c13}, Ch.~18 \chref{c14});
  \item you only need to \emph{count} and \(\bigO{\log^3 n}\) already fits -- the merge-sort tree (\(\secref{c19-mst}\)) is simpler to code and to reason about;
  \item the values sit on a \emph{tree path} \(u\text{--}v\), not a subarray -- that is \(\secref{c19-ontrees}\), four roots instead of two.
}
One-line test: \emph{a static array and a per-subarray order statistic or count -- make one frequency tree per prefix, then a range is the difference of two prefixes and the \(k\)-th is a walk down that difference.}
\end{recognitionbox}

% ---------------------------------------------------------------
\subsection{Intuition: a range is the difference of two prefixes}
% ---------------------------------------------------------------

The merge-sort tree counts by cutting the \emph{index} range into \(O(\log n)\) nodes and bisecting a \emph{value} list in each -- a genuine \(\log\times\log\). Chairman removes one \(\log\) by moving the whole story to the value axis. A frequency segment tree over values can already answer ``how many elements \(\le x\) overall'' in \(O(\log V)\). The only thing it lacks is \emph{which} elements -- a subarray \(A[l..r]\) is not ``all'' elements. But \(A[l..r] = A[1..r]\setminus A[1..l-1]\), a difference of two prefixes, and counts are additive over set difference. So if we keep a frequency tree \emph{for every prefix}, we recover any subarray's histogram by subtracting two trees.

Keeping \(n\) separate frequency trees looks like \(O(nV)\) memory -- hopeless. \(\secref{c19-copy}\) already showed the fix: consecutive prefixes \(\mathit{root}_{i-1}\) and \(\mathit{root}_i\) differ by \emph{one} inserted value, so \(\mathit{root}_i\) shares almost all of \(\mathit{root}_{i-1}\) and only path-copies the \(\bigO{\log V}\) nodes from the root down to leaf \(\mathit{rank}(A_i)\). The chain \(\mathit{root}_0\dots\mathit{root}_n\) costs \(\bigO{n\log V}\) nodes total -- the same budget as a merge-sort tree, but now every range count is \emph{one} subtraction of two node fields, and the \(k\)-th is a walk.

Why the \(k\)-th is a walk and not a search: at any value node \([p,q]\) split at \(\mathit{mid}\), the number of query elements landing in the lower half \([p,\mathit{mid}]\) is \((\mathit{root}_r.\mathit{left}.\mathit{cnt})-(\mathit{root}_{l-1}.\mathit{left}.\mathit{cnt})\) -- available with two pointer reads. If \(k\) is at most that, the answer is in the lower half; else it is in the upper half with \(k\) reduced by that count. Descending \(\mathit{root}_{l-1}\) and \(\mathit{root}_r\) \emph{together}, one decision per value level, reads the answer leaf in \(\bigO{\log V}\). There is no outer binary search over the value -- the tree descent \emph{is} the search (Ch.~19's \(\secref{c13-search}\) ``descend, do not bisect outside'' lesson, now on the value axis with a difference of two versions).

\begin{definition}{Prefix-version (chairman) tree}{c19-roots}
Let \(A[1..n]\) be a static array with distinct sorted values \(\mathit{vals}_1<\dots<\mathit{vals}_V\), and write \(\rho_i=\mathit{rank}(A_i)\in[1,V]\). A \term{chairman tree} is a chain of roots \(\mathit{root}_0,\dots,\mathit{root}_n\) of frequency segment trees over the value range \([1,V]\), built by \(\mathit{root}_0\coloneqq\) (all-zero tree) and \(\mathit{root}_i\coloneqq\texttt{insert}(\mathit{root}_{i-1},\rho_i)\), where \texttt{insert} path-copies the nodes on the root-to-leaf-\(\rho_i\) route, setting each copied node's \(\mathit{cnt}\) to its predecessor's \(\mathit{cnt}+1\) and sharing the untouched child pointer. Each node's \(\mathit{cnt}\) is the number of prefix elements whose rank lies in that node's value interval, so the range \(A[l..r]\)'s histogram is the \term{two-root difference} \(\mathit{root}_r-\mathit{root}_{l-1}\) taken node-wise.
\end{definition}

\begin{theorem}{Two-root counting and \(k\)-th selection on chairman trees}{c19-kth}
On \(n\) elements over a value range of size \(V\), the chairman chain uses \(\bigO{n\log V}\) nodes and builds in \(\bigO{n\log V}\) time. For \(1\le l\le r\le n\): the number of elements of \(A[l..r]\) whose rank is \(\le t\) is \(\mathit{root}_r\) and \(\mathit{root}_{l-1}\)'s prefix-of-value counts differ by \(\bigO{\log V}\) node reads; in particular \(\mathit{count}(\le \mathit{vals}_t\ \text{in}\ [l,r])\) is obtained by descending to value node \([1,t]\) and subtracting the two accumulated left-counts, in \(\bigO{\log V}\). The \(k\)-th smallest of \(A[l..r]\) is found by a single joint descent of \(\mathit{root}_{l-1}\) and \(\mathit{root}_r\), reading two left-child counts per level, in \(\bigO{\log V}\) time -- no outer binary search over values.
\end{theorem}
\begin{proofbox}
\textbf{Memory.} Each \texttt{insert} clones one root-to-leaf path of \(\lceil\log_2 V\rceil+1=O(\log V)\) nodes and shares the rest (\(\thmref{c19-memory}\), reused with \(n\coloneqq V\) as the tree width); \(n\) inserts give \(\bigO{n\log V}\).\\
\textbf{Count.} By construction a node of \(\mathit{root}_i\) over value interval \([p,q]\) has \(\mathit{cnt}=\#\{j\le i:\rho_j\in[p,q]\}\). Since \(A[l..r]=A[1..r]\setminus A[1..l-1]\) and this is a disjoint subtraction, the number of query ranks in \([p,q]\) is \(\mathit{root}_r.\mathit{node}_{[p,q]}.\mathit{cnt}-\mathit{root}_{l-1}.\mathit{node}_{[p,q]}.\mathit{cnt}\). The value set \(\le\mathit{vals}_t\) is a union of \(O(\log V)\) canonical value nodes; summing the two-root differences over them is one \(\bigO{\log V}\) descent (equivalently, \(\mathit{count}\) is the prefix-sum of the difference histogram).\\
\textbf{\(k\)-th.} Maintain the invariant ``\(k\) counts elements not yet placed, within the current value interval \([p,q]\).'' At \([p,q]\) with left child \([p,\mathit{mid}]\), let \(c=(\mathit{root}_r.\mathit{left}.\mathit{cnt})-(\mathit{root}_{l-1}.\mathit{left}.\mathit{cnt})\): by the Count paragraph \(c\) is exactly the number of \(A[l..r]\) elements landing in \([p,\mathit{mid}]\). If \(k\le c\) the \(k\)-th is there and \(([p,\mathit{mid}],k)\) keeps the invariant; else it is in \([\mathit{mid}+1,q]\) with \(k-c\). Each step halves the value interval and costs two \(\mathit{cnt}\) reads, so after \(\bigO{\log V}\) steps the interval is a single leaf whose value is the answer. Distinctness of \(A\) is \emph{not} required: equal values simply share a rank and their counts add.
\end{proofbox}

\noindent\textbf{In plain words.}  Number the boxes by \emph{value}, one box per distinct number (after compression). Walking the array left to right, drop each element into its value-box -- and \emph{photograph the whole shelf after every drop}. Because consecutive photos differ in one box only, each photo is stored as a tiny edit of the previous one (path copying). To ask about a subarray \(A[l..r]\), subtract photo \(l-1\) from photo \(r\): what remains is the multiset of the range. ``How many \(\le x\)'' becomes ``read one number off each photo and subtract.'' ``What is the \(k\)-th smallest'' becomes ``walk down the two shelves side by side; at every fork, the gap between the two left-halves tells you which fork holds element \(k\).''

% ---------------------------------------------------------------
\subsection{The build (one path copy per element) and the joint \(k\)-th}
% ---------------------------------------------------------------

\(\coderef{c19-chairman}\) is \(\coderef{c19-pseg}\)'s \texttt{update} specialised to a frequency count (bump \(\mathit{cnt}\), not a leaf value) and run once per array element to make the root chain. \(\coderef{c19-kth}\) walks two roots at once.

\begin{lstlisting}[caption={chairman build: node pool with a shared null at index 0; ROOT(i) = ROOT(i-1) with one path copy inserting rank(A(i)); O(log V) new nodes per element}, label={lst:c19-chairman}]
struct Node { int cnt, l, r; };
vector<Node> T;                       // node pool; T[0] = the all-zero "null" node
int newnode(int cnt, int l, int r){ T.push_back({cnt,l,r}); return (int)T.size()-1; }
int ROOT[MAXn];                       // ROOT[i] = frequency tree over A[1..i]
int V;                                // number of distinct (compressed) values

int insert(int prev, int tl, int tr, int pos){     // path-copy prev, add one at value-rank pos
    int cur = newnode(T[prev].cnt + 1, T[prev].l, T[prev].r);  // clone node, bump cnt
    if (tl != tr){
        int tm = (tl + tr) / 2;
        if (pos <= tm){
            int nl = insert(T[prev].l, tl, tm, pos);   // recurse into LEFT child
            T[cur].l = nl;                              //   new left, right stays SHARED
        } else {
            int nr = insert(T[prev].r, tm+1, tr, pos); // recurse into RIGHT child
            T[cur].r = nr;                              //   new right, left stays SHARED
        }
    }
    return cur;
}
void build(const vector<int>& rk){          // rk[i] = compressed rank of A[i], in 1..V
    T = { {0,0,0} };                        // node 0: cnt 0, both children point to itself
    ROOT[0] = 0;
    for (int i = 1; i <= (int)rk.size()-1; i++)
        ROOT[i] = insert(ROOT[i-1], 1, V, rk[i]);
}
\end{lstlisting}

\begin{lstlisting}[caption={k-th smallest of A(l..r) by descending ROOT(l-1) and ROOT(r) together -- the left-count DIFFERENCE decides each fork; O(log V)}, label={lst:c19-kth}]
// u = ROOT[l-1], v = ROOT[r]; value interval [tl..tr]; want the k-th smallest (1-based)
int kth(int u, int v, int tl, int tr, int k){
    if (tl == tr) return tl;                 // leaf rank = the answer's compressed value
    int tm = (tl + tr) / 2;
    // # query elements whose rank is in the LOWER value-half [tl..tm]:
    int leftCount = T[T[v].l].cnt - T[T[u].l].cnt;
    if (k <= leftCount)                       // it lives in [tl..tm]: descend both left children
        return kth(T[u].l, T[v].l, tl, tm, k);
    else                                      // it lives in [tm+1..tr], past leftCount elements
        return kth(T[u].r, T[v].r, tm+1, tr, k - leftCount);
}
// "how many of A[l..r] are <= vals[t]": descend to value prefix [1..t] subtracting counts.
\end{lstlisting}

Two things to hold onto. First, \(\mathit{leftCount}\) is \emph{the whole query}: one subtraction of two pool reads, and it is exactly the merge-sort tree's count restricted to the lower half -- but paid for in \(O(1)\) instead of a bisection per cover node. Second, when the array has duplicates, two equal elements map to the same rank and their counts simply add inside that leaf; the descent is unchanged (MKTHNUM happens to promise distinct values, but nothing here needs it).

% ---------------------------------------------------------------
\subsection{The hand trace on \(A\)}
% ---------------------------------------------------------------

Take the value axis as the eight distinct values \(\{0,1,2,3,4,5,6,7\}\) (already compressed, in order). Let \(F_i\) be the frequency histogram of \(A[1..i]\) over values \(0,1,2,3,4,5,6,7\) -- i.e.\ the leaf vector of \(\mathit{root}_i\). A few prefixes of the chain, read straight off \(A=(3,2,4,5,1,1,5,3,7,0,2,4,1,6,3,2)\):
\[
\begin{array}{c|cccccccc|c}
\text{prefix } i & 0&1&2&3&4&5&6&7 & A[1..i]\ \text{adds}\\\hline
\mathit{root}_0=F_0 & 0&0&0&0&0&0&0&0 & (\varnothing)\\
\mathit{root}_2=F_2 & 0&0&1&1&0&0&0&0 & A_2=2\\
\mathit{root}_8=F_8 & 0&2&1&2&1&2&0&0 & A_8=3\\
\mathit{root}_9=F_9 & 0&2&1&2&1&2&0&1 & A_9=7
\end{array}
\]
Each row is the previous row with \(1\) added in the column of \(A_i\); \(\mathit{root}_i\) is the path copy of \(\mathit{root}_{i-1}\) that walks to leaf \(A_i\) and bumps \(\mathit{cnt}\) on those \(\lceil\log_2 8\rceil+1=4\) nodes.

\begin{figure}[H]
\centering
\begin{tikzpicture}[diagram, scale=0.62, transform shape]
  % two roots ROOT_2 and ROOT_9, joint descent for kth[3..9], k=4
  \node[boxnode, primaryblue, minimum size=10mm, font=\scriptsize] (r9) at (0,2.6) {\(\mathit{root}_9\)};
  \node[boxnode, primaryblue, minimum size=10mm, font=\scriptsize] (r2) at (4.4,2.6) {\(\mathit{root}_2\)};
  \node[font=\tiny, text=darkgray] at (2.2,2.6) {two roots stand side by side};
  % level 1: value [0..7] split into [0..3] / [4..7]; leftCount=3
  \node[boxnode, difforange, minimum size=9mm, font=\scriptsize] (L) at (1.1,1.2) {\([0..3]\)};
  \node[boxnode, gray, minimum size=9mm, font=\scriptsize] (R) at (3.3,1.2) {\([4..7]\)};
  \node[font=\tiny\ttfamily, text=orange!60!black, anchor=north] at (1.1,0.75) {\(c_L{=}3{<}4\)};
  \node[font=\tiny\ttfamily, text=darkgray, anchor=north] at (3.3,0.75) {\(k'{=}4{-}3{=}1\)};
  % level 2 within [4..7]: [4..5] / [6..7]
  \node[boxnode, difforange, minimum size=9mm, font=\scriptsize] (LL) at (2.3,-0.2) {\([4..5]\)};
  \node[boxnode, gray, minimum size=9mm, font=\scriptsize] (RR) at (4.2,-0.2) {\([6..7]\)};
  \node[font=\tiny\ttfamily, text=orange!60!black, anchor=north] at (2.3,-0.65) {\(c{=}3{\ge}1\)};
  % level 3 within [4..5]: pick 4
  \node[boxnode, successgreen, minimum size=8mm, font=\scriptsize] (four) at (3.1,-1.5) {\(\{4\}\)};
  \node[font=\tiny\ttfamily, text=green!45!black, anchor=north] at (3.1,-1.95) {answer 4};
  \draw[flow] (r9)--(L); \draw[flow] (r2)--(L);
  \draw[flow] (r9)--(R); \draw[flow] (r2)--(R);
  \draw[flow] (R)--(LL); \draw[flow] (R)--(RR);
  \draw[flow] (LL)--(four);
\end{tikzpicture}
\caption{The joint two-root descent for the \(k\)-th smallest of \(A[3..9]\) (\(k=4\)).  At each value node the branch label is the \emph{difference} of \(\mathit{root}_9\)'s and \(\mathit{root}_2\)'s left-child counts -- the number of query elements in the lower half.  \(3<4\) pushes into \([4..7]\) with \(k'=1\); then \([4..5]\) holds \(3\ge1\) elements so we take the low half and land on value \(4\).  Orange = descended, gray = skipped.}
\label{fig:c19-roots}
\end{figure}

\begin{tracebox}[kth on (3..9) with k equal 4 gives 4 -- subtracting version 9 from version 2]
\textbf{Range histogram.} \(F_9-F_2=(0,2,1,2,1,2,0,1)-(0,0,1,1,0,0,0,0)=(0,2,0,1,1,2,0,1)\) over values \(0..7\): the multiset \(A[3..9]=\{1,1,3,4,5,5,7\}\) (length 7, and the counts sum to 7).\\
\textbf{Fork 1, node \([0..7]\), left half \([0..3]\):} leftCount \(=(0+2+1+2)-(0+0+1+1)=5-2=3\).  \(k=4>3\) \(\Rightarrow\) go to \([4..7]\) with \(k'=4-3=1\).\\
\textbf{Fork 2, node \([4..7]\), left half \([4..5]\):} \(F\) difference there \(=(1+2)-0=3\).  \(k'=1\le3\) \(\Rightarrow\) go to \([4..5]\), still \(k'=1\).\\
\textbf{Fork 3, node \([4..5]\), left half \([4]\):} count \(=1-0=1\).  \(k'=1\le1\) \(\Rightarrow\) value \(\boxed{4}\).\\
\textbf{Cross-check} by sorting: \(A[3..9]=(4,5,1,1,5,3,7)\to(1,1,3,\underline{4},5,5,7)\); the 4th is \(4\).  Three pointer-pairs, no re-sorting and no outer binary search.
\end{tracebox}

\begin{tracebox}[kth(1..8, k=5) = 3, a second two-root walk on the same chain]
\textbf{Range histogram} \(F_8-F_0=(0,2,1,2,1,2,0,0)\): multiset \(\{1,1,2,3,3,4,5,5\}\).\\
\textbf{Fork 1 \([0..7]\), left \([0..3]\):} count \(=0+2+1+2=5\).  \(k=5\le5\) \(\Rightarrow\) stay left \([0..3]\).\\
\textbf{Fork 2 \([0..3]\), left \([0..1]\):} count \(=0+2=2\).  \(5>2\) \(\Rightarrow\) go \([2..3]\) with \(k'=5-2=3\).\\
\textbf{Fork 3 \([2..3]\), left \([2]\):} count \(=1\).  \(3>1\) \(\Rightarrow\) go \([3]\) with \(k''=3-1=2\), leaf \(=\) value \(\boxed{3\ }\) (its own rank; the two copies of \(3\) both live in this leaf).\\
\textbf{Cross-check:} sort \(A[1..8]=(3,2,4,5,1,1,5,3)\to(1,1,2,3,\underline{3},4,5,5)\); the 5th is \(3\).  The subtraction of \(\mathit{root}_0\) is all-zero, so this is just \(\mathit{root}_8\) read alone -- as a \emph{prefix} query it should be.
\end{tracebox}

\begin{complexitybox}[The chairman tree]
\textbf{Build} \(\bigO{n\log V}\) nodes and time (\(n\) inserts, one path copy each). \textbf{Range \(k\)-th} \(\bigO{\log V}\) per query (a joint descent, no outer search). \textbf{Range count} \(\le x\) \(\bigO{\log V}\). \textbf{Memory} \(\bigO{n\log V}\): with \(n\le10^5,V\le10^5\) budget \(\sim n\cdot 17\approx1.7\cdot10^6\) nodes \(\times\,3\) ints \(\approx 20\) MB -- the pool \emph{must} be sized for this, exactly as \(\thmref{c19-memory}\) warns. \textbf{Static only}, like the merge-sort tree: a write invalidates every later prefix root.
\end{complexitybox}

% ---------------------------------------------------------------
\subsection{Problem Lab: SPOJ MKTHNUM -- the \(k\)-th number in a subarray}
% ---------------------------------------------------------------

\solvedproblem{c19-mkthnum}{K-th Number}{SPOJ MKTHNUM}
  {https://www.spoj.com/problems/MKTHNUM/}{SPOJ L6}
  {the chairman tree in one move: prefix-frequency roots, answer each query by a two-root joint descent}

\begin{problemstatement}[SPOJ MKTHNUM K-th Number]
Given an array \(a[1..n]\) of \(n\) \emph{distinct} integers, answer \(m\) queries \(Q(i,j,k)\): the \(k\)-th number in \(a[i..j]\) if that segment were sorted in non-decreasing order.

\textbf{Input.} Line 1: \(n\) and \(m\) (\(1\le n\le 100000\), \(1\le m\le 5000\)). Line 2: \(n\) distinct integers, \(|a_i|\le 10^9\). Then \(m\) lines, each \(i\ j\ k\) with \(1\le i\le j\le n\), \(1\le k\le j-i+1\).

\textbf{Output.} For each query, the \(k\)-th smallest of \(a[i..j]\), one per line.

\textbf{Example (official).}
\begin{verbatim}
Input:
7 3
1 5 2 6 3 7 4
2 5 3
4 4 1
1 7 3

Output:
5
6
3
\end{verbatim}
\end{problemstatement}

\begin{attemptgate}[Attempt it yourself first -- \lv{4}, give yourself 30 minutes]
Sorting every query's subarray is \(\bigO{m\cdot n\log n}\) with \(n=10^5\): a \tle. A merge-sort tree answers \(k\)-th in \(\bigO{\log^3 n}\approx 17^3\cdot 5000\approx2.4\cdot10^7\) -- passable, but this is the canonical chairman problem. The right frame: the array is fixed, so precompute one \emph{prefix} frequency tree per position, then each query is a two-root walk.
\peek{which two versions of a growing histogram subtract to give \(a[i..j]\)?}
\nudge{coordinate-compress \(a\) to ranks \(1..n\); build \(\mathit{root}_t=\mathit{insert}(\mathit{root}_{t-1},\mathit{rank}(a_t))\); the \(k\)-th of \(a[i..j]\) is a joint descent of \(\mathit{root}_{i-1}\) and \(\mathit{root}_{j}\).}
\fullhint{Compress to ranks, build the chain (\(\coderef{c19-chairman}\)). For each query call \(t\leftarrow\mathit{kth}(\mathit{root}_{i-1},\mathit{root}_{j},1,n,k)\) (\(\coderef{c19-kth}\)) and print \(\mathit{vals}[t]\). Total \(\bigO{(n+m)\log n}\approx(10^5+5\cdot10^3)\cdot17\approx1.8\cdot10^6\) node steps -- a comfortable \ac. Use the pooled \(\mathit{cnt}\) as an \code{int}; \(\mathit{vals}\) can be \code{int} (\(|a_i|\le10^9\)).}
\end{attemptgate}

\solstep{Understanding}
\(n\le10^5\) distinct values, \(m\le5000\) range \(k\)-th queries. Because the array never changes, a prefix-frequency-tree chain answers any subarray histogram by subtraction. This is chairman's textbook use: each query is one joint descent, \(\bigO{\log n}\), with values mapped to ranks so negative or \(10^9\)-sized inputs are no issue.

\solstep{Key Observations}
\begin{observation}
\(a[i..j]\) is \(A[1..j]\) minus \(A[1..i-1]\), so \emph{every} range statistic that is a sum of per-value indicators becomes a difference of two prefix versions -- the whole reason the versions are prefixes and not arbitrary.
\end{observation}
\begin{observation}
The \(k\)-th is a \emph{descent}, not a search: at each value level the left-count difference is the exact number of query elements below the midpoint, so it decides the branch outright (\(\secref{c13-search}\)'s ``descend, don't bisect outside,'' applied to the value axis).
\end{observation}
\begin{observation}
Distinctness is an \emph{input promise} of MKTHNUM, not a requirement of the structure -- duplicates land in the same leaf and add; \(\coderef{c19-kth}\) needs no change.
\end{observation}

\solstep{Think It Through With Me}
\begin{thinkbox}
\deadend{Sort each subarray, or keep one merge-sort tree}
{Sorting is \(\bigO{m\,n\log n}\) -- \tle. The merge-sort tree gets \(k\)-th via an outer value binary search at \(\bigO{\log^3 n}\): correct, but you pay a \(\log\) more than needed and re-tile the index range at every step.}
\deadend{A Fenwick tree over the values, sweeping}
{Works \emph{offline} for a count \(\le x\) if queries don't need order statistics; but the \(k\)-th is not a linear query -- you'd binary-search the answer against the BIT, \(\bigO{\log^2 n}\), and it must be re-run per query. The prefix-version chain answers \emph{online} and in one descent.}
\aha{Make a frequency segment tree \emph{for each prefix}, cheaply by path copying, and the \(k\)-th in \(a[i..j]\) is a single \(\bigO{\log n}\) walk down \(\mathit{root}_j\) and \(\mathit{root}_{i-1}\) side by side.}
\end{thinkbox}

\solstep{The Approach Ladder}
\begin{center}\renewcommand{\arraystretch}{1.25}\setlength{\tabcolsep}{3pt}
\begin{tabular}{@{}c p{0.22\linewidth} p{0.26\linewidth} l c@{}}
\toprule
\textbf{\#} & \textbf{Approach} & \textbf{Idea} & \textbf{Time} & \textbf{Verdict}\\
\midrule
1 & \lv{2}\ sort per query & copy \(a[i..j]\), sort, index \(k\) & \(\bigO{m\,n\log n}\) & \tle \\
2 & \lv{3}\ merge-sort tree & value-bisect an outer binary search (\(\secref{c19-mst}\)) & \(\bigO{m\log^3 n}\) & \ac \\
3 & \lv{4}\ offline BIT + bsearch & sweep by threshold, search the answer & \(\bigO{m\log^2 n}\) & \ac \\
4 & \lv{4}\ \textbf{chairman tree} & prefix versions; joint two-root descent (\(\coderef{c19-kth}\)) & \(\bigO{(n+m)\log n}\) & \ac \\
\bottomrule
\end{tabular}\end{center}
Rungs~2 and~3 are correct but re-pay a logarithm; the chairman tree is rung~4 -- one path copy per element at build, one clean descent per query. That \(\bigO{\log n}\) per query with no outer loop is the whole reason this structure exists.

\solstep{Solution}
Coordinate-compress \(a\) to ranks. Build the prefix chain \(\mathit{root}_0..\mathit{root}_n\) (\(\coderef{c19-chairman}\)). For each query \((i,j,k)\), run the joint descent \(t=\mathit{kth}(\mathit{root}_{i-1},\mathit{root}_{j},1,V,k)\) (\(\coderef{c19-kth}\)) and print the original value \(\mathit{vals}[t]\).

\solstep{Proof of Correctness}
\begin{proofbox}
\textbf{Chain is right.} \(\mathit{root}_0\) is the empty histogram; \(\mathit{root}_t\) path-copies \(\mathit{root}_{t-1}\) and adds \(1\) along the single root-to-leaf path to \(\mathit{rank}(a_t)\), so by induction each node's \(\mathit{cnt}\) equals the number of \(a_1..a_t\) with rank in that node's value interval, and every other node is \emph{shared} (never mutated) -- exactly \(\thmref{c19-memory}\)'s invariant.\\
\textbf{Two-root answer is right.} By \thmref{c19-kth}, at each descent the quantity \(T[\mathit{root}_j.\mathit{left}].\mathit{cnt}-T[\mathit{root}_{i-1}.\mathit{left}].\mathit{cnt}\) is the number of elements of \(a[i..j]\) in the lower value half; branching on whether \(k\) exceeds it keeps the invariant ``the answer's rank is the \(k\)-th smallest among the still-live query elements within the current interval,'' shrinking to the answer leaf in \(\bigO{\log V}\). Printing \(\mathit{vals}[t]\) un-compresses the rank.\\
\textbf{Complexity.} Build \(\bigO{n\log n}\); query \(\bigO{\log n}\). With \(n=10^5,m\le5\cdot10^3\) that is \(\approx1.8\cdot10^6\) node steps, \ac; node count \(\approx n(\lceil\log_2 n\rceil+1)+1\approx1.8\cdot10^6\), within SPOJ memory when the pool is sized for it.
\end{proofbox}

\solstep{Implementation}
\begin{lstlisting}[caption={SPOJ MKTHNUM full solution: compress, build the prefix chain, answer each query with one joint two-root descent}, label={lst:c19-mkthnum-soln}]
#include <bits/stdc++.h>
using namespace std;
const int MAXn = 100005;
struct Node { int cnt, l, r; };
vector<Node> T;
int newnode(int cnt,int l,int r){ T.push_back({cnt,l,r}); return (int)T.size()-1; }
int ROOT[MAXn], rk[MAXn];
vector<int> vals;               // sorted distinct values (rank r -> vals[r-1])
int V;

int insert(int prev, int tl, int tr, int pos){
    int cur = newnode(T[prev].cnt + 1, T[prev].l, T[prev].r);
    if (tl != tr){
        int tm = (tl + tr) / 2;
        if (pos <= tm){ int nl = insert(T[prev].l, tl, tm, pos); T[cur].l = nl; }
        else          { int nr = insert(T[prev].r, tm+1, tr, pos); T[cur].r = nr; }
    }
    return cur;
}
int kth(int u, int v, int tl, int tr, int k){       // u=ROOT[i-1], v=ROOT[j]
    if (tl == tr) return tl;
    int tm = (tl + tr) / 2;
    int leftCount = T[T[v].l].cnt - T[T[u].l].cnt;
    if (k <= leftCount) return kth(T[u].l, T[v].l, tl, tm, k);
    else                return kth(T[u].r, T[v].r, tm+1, tr, k - leftCount);
}
int main(){
    int n, m; scanf("%d%d", &n, &m);
    vector<int> A(n+1);
    vals.reserve(n);
    for (int i=1;i<=n;i++){ scanf("%d",&A[i]); vals.push_back(A[i]); }
    sort(vals.begin(), vals.end());
    vals.erase(unique(vals.begin(), vals.end()), vals.end());
    V = vals.size();
    for (int i=1;i<=n;i++)
        rk[i] = lower_bound(vals.begin(), vals.end(), A[i]) - vals.begin() + 1;
    T = { {0,0,0} }; ROOT[0] = 0;
    for (int i=1;i<=n;i++) ROOT[i] = insert(ROOT[i-1], 1, V, rk[i]);
    while (m--){
        int i, j, k; scanf("%d%d%d", &i, &j, &k);
        int t = kth(ROOT[i-1], ROOT[j], 1, V, k);
        printf("%d\n", vals[t-1]);
    }
}
\end{lstlisting}

\solstep{Trace on the Official Sample}
\begin{tracebox}[MKTHNUM on a=(1,5,2,6,3,7,4): three two-root walks]
Values already distinct, sorted \(\mathit{vals}=(1,2,3,4,5,6,7)\), ranks \(=\) values here.\\
\textbf{(2,5,3)}: \(a[2..5]=(5,2,6,3)\to\) sorted \((2,3,5,6)\). \(\mathit{kth}(\mathit{root}_1,\mathit{root}_5,3\)\ drives to rank~5, \(\mathit{vals}[5]=\boxed{5}\).\\
\textbf{(4,4,1)}: \(a[4..4]=(6)\). Only one element, so the descent lands on rank~6, \(\mathit{vals}[6]=\boxed{6}\).\\
\textbf{(1,7,3)}: \(a=(1,5,2,6,3,7,4)\to(1,2,3,4,5,6,7)\); \(\mathit{root}_0\) is all-zero, so this reads \(\mathit{root}_7\) alone -- 3rd smallest \(\to\mathit{vals}[3]=\boxed{3}\).  Matches official output \(5,\;6,\;3\).
\end{tracebox}

\solstep{Complexity}
\begin{complexitybox}
Compress \(\bigO{n\log n}\); build \(\bigO{n\log n}\approx1.8\cdot10^6\) node touches; queries \(\bigO{m\log n}\approx5\cdot10^3\cdot17\approx8.5\cdot10^4\) descents -- all \ac. Memory \(\approx1.8\cdot10^6\) nodes of \(3\) ints \(\approx21\) MB plus the \(\mathit{vals}/\mathit{ROOT}/\mathit{rk}\) arrays -- size the pool for \(\bigO{n\log n}\), never \(4n\).
\end{complexitybox}

\solstep{Verification and Edge Cases}
\begin{center}\renewcommand{\arraystretch}{1.2}
\begin{tabular}{@{}p{0.22\linewidth} p{0.30\linewidth} p{0.36\linewidth}@{}}
\toprule
\textbf{Case} & \textbf{Why tricky} & \textbf{Code behaviour}\\
\midrule
\(i=j\) single slot & range length 1 & \(\mathit{kth}\) returns that rank; \(k=1\) forced \\
\(k=1\) / \(k=j-i+1\) & extreme descent & always takes the left or right edge -- \(\bigO{\log V}\)\\
negative \(a_i\) & rank mapping & compression is over sorted \(\mathit{vals}\); sign irrelevant \\
\(|a_i|=10^9\) & int ceiling & ranks \(\le n\le10^5\) fit \code{int}; store values \code{int} \\
all-equal prefix & shared child & path copy shares the untouched subtree; \(\mathit{cnt}\) still exact \\
\bottomrule
\end{tabular}\end{center}

\begin{hardeningbox}
\begin{itemize}[nosep,leftmargin=1.4em]
  \item \textbf{Pool size is the number-one MLE.} \(\sim n(\log_2 n+1)\) nodes, not \(4n\). On \(n=10^5\) that is about \(1.8\cdot10^6\); forget the multiplier and the judge kills you at the largest test.
  \item \textbf{Subtract the \emph{right} pair of roots.} The range is \(\mathit{root}_j-\mathit{root}_{i-1}\); using \(\mathit{root}_{i}\) drops \(a_i\) from the histogram (the classic off-by-one that silently loses one element).
  \item \textbf{Read children of the \emph{pool}, not of \(\mathit{root}\).} The descent uses \(T[T[v].l].\mathit{cnt}-T[T[u].l].\mathit{cnt}\) -- two indirections; reading \(T[v].\mathit{cnt}\) directly compares whole-tree counts and cannot choose a half.
  \item \textbf{Compress before you build.} Building over raw \([-10^9,10^9]\) needs a balanced-over-values tree of depth \(30\) and \(O(\log 10^9)\) memory per insert -- compression to \(1..V\) makes the depth \(\log n\).
  \item \textbf{Un-compress the answer.} \(\mathit{kth}\) returns a rank; print \(\mathit{vals}[\text{rank}-1]\), not the rank.
\end{itemize}
\end{hardeningbox}

% ---------------------------------------------------------------
\subsection{Common Pitfalls}
% ---------------------------------------------------------------

\begin{warningbox}[Trap 1: sizing the pool like a mutable tree]
A single heap tree fits in \(4n\). A chairman \emph{chain} is \(n\) inserts each cloning \(\log V\) nodes: budget \(\bigO{n\log n}\). Under-sizing is an out-of-bounds write into the pool -- a wrong answer or crash that looks unrelated to the algorithm.
\end{warningbox}

\begin{warningbox}[Trap 2: the off-by-one \(\mathit{root}_{l-1}\)]
The prefix-version range is `everything up to \(r\) minus everything up to \(l-1\)'. Reaching for \(\mathit{root}_l\) (not \(\mathit{root}_{l-1}\)) silently deletes \(a_l\) from the multiset; every count and the \(k\)-th shift. Draw \(A[1..r]\setminus A[1..l-1]=A[l..r]\) when in doubt.
\end{warningbox}

\begin{warningbox}[Trap 3: mutating a shared child]
\(\mathit{root}_i\) shares most pointers with \(\mathit{root}_{i-1}\). \emph{Never} overwrite \(T[\mathit{root}_{i-1}].\mathit{cnt}\) in place -- you corrupt every version downstream of it. \(\coderef{c19-chairman}\) avoids this by cloning the node first (\code{newnode}) and only writing the clone. Persistence's one law: allocate, then fill; never edit a borrowed node.
\end{warningbox}

\begin{warningbox}[Trap 4: comparing wrong in the descent]
The branch test is \(k\le\text{leftCount}\) where \texttt{leftCount} is a \emph{difference of left-child counts}. Using the whole-node count, or forgetting to reduce \(k\) by \texttt{leftCount} when going right, mis-places the answer by a whole value-half.
\end{warningbox}

% ---------------------------------------------------------------
\subsection{A re-treat in prose: CF 813E ``Army Creation''}
% ---------------------------------------------------------------

\begin{insightbox}[The chairman twin of Ch.~19's sliding window -- same task, other engine]
Chapter~19 (\chref{c13}) solved \emph{CF~813E Army Creation} \emph{offline}: reading the array left to right, keep a BIT/segment tree that, at each index \(p\), stores \(1\) iff \(p\) is the \emph{most recent} occurrence of its value; a query \([l,r]\) then asks for the number of such markers in \([l,r]\), i.e.\ the number of \emph{distinct} values in the range -- a sweep sorted by right endpoint. That is a fine, fast solution. Now notice what it was hiding: it maintained a per-prefix `latest-occurrence set' and only ever \emph{read} one version (the current one) at a time. \emph{Persisting} those prefixes is exactly a chairman-style chain: version \(t\) is the tree over \emph{positions} where value \(a_t\)'s marker has been moved from \(\mathit{prev}[t]\) to \(t\) -- two point-updates on \(\mathit{root}_{t-1}\)'s copy, so \(\mathit{root}_t\) costs \(\bigO{\log n}\) new nodes. Then a query ``distinct values in \(A[l..r]\)'' is a \emph{range sum on a single stored version}: \(\mathit{root}_r\) already knows which of \(A[1..r]\)'s values have their last occurrence \(\ge l\), so counting markers in \([l,r]\) of version \(\mathit{root}_r\) is \(\bigO{\log n}\) -- \emph{online}, in any query order, no offline sort. Same law as \(\coderef{c19-chairman}\), applied to `latest occurrence' instead of `frequency'. (We re-treat the idea only; the problem's URL stays with Ch.~19's solution.)
\end{insightbox}

This is the Ch.~18/Ch.~20 handover pattern in miniature: the offline sweep and the persistent chain solve the same range-distinct task, and the second trades a sort for \(\bigO{\log n}\)-memory-per-step so that queries can interleave. When you see ``answer about \(A[l..r]\) in arbitrary order,'' ask whether a prefix version already holds the answer in one root.

% ---------------------------------------------------------------
\subsection{Practice Problems}
% ---------------------------------------------------------------

\begin{stuckbox}
\begin{itemize}[nosep,leftmargin=1.4em]
  \item Did you \emph{coordinate-compress} before building, so the tree depth is \(\log n\) and not \(\log 10^9\)?
  \item Is the pool sized for \(\bigO{n\log n}\) nodes (\(\sim n(\log_2 n+1)\)), not \(4n\)?
  \item Does the range use \(\mathit{root}_r-\mathit{root}_{l-1}\) (not \(\mathit{root}_l\))?
  \item In the descent, is \texttt{leftCount} a difference of \emph{left-child} counts, and do you subtract it from \(k\) when going right?
  \item Do you print \(\mathit{vals}[\text{rank}-1]\), un-compressing the leaf you landed on?
\end{itemize}
\end{stuckbox}

\begin{practiceset}
  \practice{Greetings}{CF 1915F}
    {https://codeforces.com/problemset/problem/1915/F}
    {CF 1500 L3}
    {You merge \(k\) lists into one sequence and must count \emph{pairs} of equal values \((i<j)\) whose original list indices are inverted (\(i\)'s list number exceeds \(j\)'s). The engine: for each \emph{value}, keep the increasing list of its positions grouped by source list, then count, for that value, the pairs where a later-list occurrence precedes an earlier-list one -- i.e.\ an inversion count \emph{within a single value's ordered position list}. That is the value-keyed `sorted positions' view of this section, once per distinct value; total the per-value answers. Solve the small case by two nested loops, then speed the `how many of this value's occurrences sit before position \(p\)' step with a binary search on the value's sorted list (a merge-sort-tree-flavoured \(\coderef{c19-mst-count}\)) or a Fenwick over the source-list numbers (Ch.~18, \chref{c14}). Rating \lv{3}; the reduction to `count occurrences of a value in a range' is exactly chairman's muscle.}
  \practice{Mike and Geometry Problem}{CF 689E}
    {https://codeforces.com/problemset/problem/689/E}
    {CF 2000 L4}
    {Given \(n\) integer segments \([l_i,r_i]\) and \(m\) points, for each point \(p\) count the triples of segments that \emph{all contain} \(p\); sum over points. Decompose: the number of segments containing \(p\) is a \emph{prefix difference} -- \(\bigl(\#\{l_i\le p\}\bigr)-\bigl(\#\{r_i<p\}\bigr)\), two sorted-array bisections or two growing counters as \(p\) sweeps. Then the per-point answer is \(\binom{c_p}{3}\), a function of that single count. The `left-ends minus right-ends up to \(p\)' structure is the \emph{two-root subtraction} of \(\secref{c19-chairman}\) wearing an offline hat: sort the endpoints, and `how many \(l_i\le p\)' is a prefix histogram you can also serve with a persistent root if points arrive out of order. Solve with the sweep (\(\bigO{(n+m)\log n}\)), then re-solve it as a chairman count to feel the equivalence. Rating \lv{4}.}
\end{practiceset}

\paragraph{Back-cites and forwards.}
This section is path copying (\(\secref{c19-copy}\), \(\defref{c19-pathcopy}\)) run on the value axis with the versions chosen for you, so \(\coderef{c19-pseg}\) is \(\coderef{c19-chairman}\) with \(\mathit{pos}\) on a rank. The `descend, don't bisect outside' lesson is Ch.~19's \(\secref{c13-search}\); the offline-BIT alternative you are told to compare against is Ch.~18 (\chref{c14}). Forward: \(\secref{c19-ontrees}\) plants a root \emph{per node of a tree} and combines \emph{four} roots to reach a path \(u\text{--}v\) -- a direct generalisation of the two-root walk here. The merge-sort tree (\(\secref{c19-mst}\)) is the same queries at \(\bigO{\log^2 n}\)/\(\bigO{\log^3 n}\); keep both, since each is easier to code in different situations.

% ---------------------------------------------------------------
\begin{recallbox}[Section 21.3 in 10 seconds]
A \emph{chairman tree} is a chain \(\mathit{root}_0..\mathit{root}_n\) of frequency segment trees over the \emph{value} axis (compressed to \(1..V\)), where \(\mathit{root}_i=\) path-copy of \(\mathit{root}_{i-1}\) with one added element -- \(\bigO{n\log V}\) nodes, build time. Any subarray is a \emph{two-root difference}: \(\mathit{root}_r-\mathit{root}_{l-1}\). \(\mathit{count}(\le x\ \text{in}\ [l,r])\) and the \(k\)-th smallest both fall out in \(\bigO{\log V}\): descend \(\mathit{root}_{l-1}\) and \(\mathit{root}_r\) together, and at each level the left-child count \emph{difference} decides the branch -- no outer binary search (unlike the merge-sort tree's \(\bigO{\log^3 n}\)). Static only. Worked: SPOJ~MKTHNUM (\(\probref{c19-mkthnum}\), \(\coderef{c19-mkthnum-soln}\)). Re-treated in prose: CF~813E range-distinct as a persisted latest-occurrence chain.
\selftest{
  \item Why is the range \(A[l..r]\)'s histogram the difference \(\mathit{root}_r-\mathit{root}_{l-1}\)?
  \item State the node/memory budget of the chain and why \(4n\) fails.
  \item Give the three branch decisions (and their left-count differences) of \(\mathit{kth}[3..9]\), \(k=4\), ending at value \(4\).
  \item Why does the chairman \(k\)-th beat the merge-sort tree's \(\bigO{\log^3 n}\)? Name the structural reason.
  \item What does coordinate compression change, and why is it needed before the build?
}
\end{recallbox}
